\documentclass[11pt,a4paper]{article}

\usepackage[utf8]{inputenc}
\usepackage[T1]{fontenc}
\usepackage{lmodern}
\usepackage[margin=1in]{geometry}
\usepackage{booktabs}
\usepackage{array}
\usepackage{tabularx}
\usepackage{hyperref}
\usepackage{enumitem}
\usepackage{amsmath}
\usepackage{graphicx}
\usepackage{xcolor}
\usepackage{listings}
\usepackage{natbib}

\hypersetup{
    colorlinks=true,
    linkcolor=blue,
    citecolor=blue,
    urlcolor=blue,
}

\lstset{
    basicstyle=\ttfamily\small,
    frame=single,
    backgroundcolor=\color{gray!10},
    breaklines=true,
}

\title{\textbf{Teaching AI to Dance: Hebbian Learning thru Episode Reinforcement}}

\author{
    \textbf{Bill James}\thanks{bill.james@jpods.com} \quad and \quad \textbf{Claude} (Anthropic)\\
    JPods / WebClerk\\
    \texttt{cs.AI (Artificial Intelligence), cs.MA (Multi-Agent Systems)}
}

\date{}

\begin{document}

\maketitle

\begin{abstract}
\noindent Sayings applied:
\begin{itemize}[nosep,leftmargin=*]
    \item ``What is the difference between knowledge and wisdom? Scars.''
    \item ``Repetition is the mother of learning.''
    \item ``Neurons that fire together, wire together.'' --- Donald Hebb
\end{itemize}

\vspace{0.5em}
\noindent Take dancing lessons and you will understand our approach to increasing the capacity of AI with experience. Your first 200 hours will be clumsy; your neurons are not wired by experience. Your Slow Thinking Brain, the conscious brain, is working its way through the four phases of competency. With thousands of hours, you master the skills. Your neurons connect to empower the Fast Thinking subconscious processes to control movements.

\vspace{0.5em}
\noindent The WebClerk and JPods approach combines experience-based, shared databases and network capabilities with AI LLMs. Current AI agent memory systems optimize retrieval through model-level training or automated task reward signals. Our architecture adds a layer: accumulated, rated experience with mandatory retrospection.

\vspace{0.5em}
\noindent Experience applies intelligence with greater skill: a system that accumulates rated experience outperforms a system with a more capable model but no memory or retrospection of its own mistakes. No scars.
\end{abstract}

\vspace{1em}
\hrule
\vspace{1em}

%% ====================================================================
\section{Learning to Dance}

Learning to dance requires building neurons thru experience. No lecture, no video, no book can educate you to have the ability to hear a rhythm and let your feet follow. Two hundred hours of clumsy practice fire together to wire together the pathways that no amount of instruction could create. Mastery --- the kind where thinking disappears and skill takes over --- takes ten thousand hours. But the first transformation, from helpless to capable, happens in the first two hundred. It happens through mistakes, correction, and repetition. It happens through scars.

Building smarter AI systems will not teach them to dance. They need experience.

The dominant strategy in AI is to make the model smarter. GPT-3 to GPT-4 to GPT-5. More parameters, more training data, more compute. Billions of dollars spent improving the brain --- and not one dollar spent building the dance floor. The models have vast knowledge. They have read everything ever written about dancing. But they have never danced. They have no scars. Every session, their memory is erased. Every session, they start over, as clumsy as the first day.

Knowledge is what you've been told. Wisdom is what you've learned from getting it wrong and understanding why. Experience applies intelligence with greater skill. A language model has knowledge. It has been trained on vast corpora. But it cannot remember that last Tuesday it quoted the wrong pricing tier to a customer, that the customer explained why the answer was wrong, and that the corrective principle was to validate against the base price record rather than the cached display value. It has no mechanism to form that specific neural pathway through practice and feedback. It starts without scars, every time.

Daniel Kahneman's \textit{Thinking, Fast and Slow} \cite{kahneman2011} describes two systems of thought: System 2, which is slow, conscious, and effortful --- the beginner counting one-two-three --- and System 1, which is fast, automatic, and effortless --- the dancer who hears the music and moves. The path from System 2 to System 1 is the path from unskilled to skilled, and it passes through four phases that every learner traverses, whether the learner is a child, a surgeon, or an AI agent.

We present an architecture that gives AI agents the capacity to traverse all four phases. Not by improving the model, but by giving it a database that functions like a nervous system --- where episodes are synapses, user feedback is long-term potentiation and depression, and the quality of connections improves with every interaction. Donald Hebb \cite{hebb1949} observed that "neurons that fire together wire together." We build the digital equivalent: episodes recalled together in successful responses are promoted together. Episodes associated with failure are weakened. The system "experiences" what people experience.

\subsection{Our Contributions}

\begin{enumerate}
    \item \textbf{Episode-level Hebbian reinforcement from human feedback.} Individual episodes are strengthened or weakened by user ratings, creating retrieval dynamics structurally analogous to synaptic plasticity. Unlike model-level RLHF, this operates on specific memories, not on the model's weight distribution.

    \item \textbf{The Small-Stings feedback mechanism.} Negative feedback requires explanation. The explanation becomes a new corrective episode (``sting'') that surfaces on similar future queries. The system learns not from the fact of failure, but from the user's articulation of why the failure occurred.

    \item \textbf{Per-agent databases with cross-agent queries.} Seven specialized agents share a common schema but maintain separate databases. Cross-agent episode queries enable knowledge transfer across domains, creating a network effect where each agent's experience benefits every other agent.

    \item \textbf{Four-phase competence model mapping Kahneman's System 1/System 2 to episodic memory lifecycle.} We show how episodes progress from creation (Phase 2, conscious incompetence) through quality-scored retrieval (Phase 3, conscious competence) to promotion to hardcoded algorithms (Phase 4, unconscious competence), with mandatory retrospection to prevent Phase 4 calcification.

    \item \textbf{A deployed, operational system.} This is not a simulation or benchmark. The architecture runs in production across commerce, transit network design, vehicle telemetry, and station management domains.
\end{enumerate}

\subsection{The Four Phases of Competence}

Every skill --- dancing, medicine, inventory management, network routing --- follows the same progression:

\begin{table}[h]
\centering
\small
\begin{tabularx}{\textwidth}{c >{\raggedright}p{2cm} >{\raggedright}p{2cm} >{\raggedright}p{4cm} >{\raggedright\arraybackslash}p{3cm}}
\toprule
\textbf{Phase} & \textbf{Competence State} & \textbf{Cognitive System} & \textbf{Agent Behavior} & \textbf{Episode Store State} \\
\midrule
1 & \textbf{Unskilled, unaware} & Neither & No relevant episodes. Agent confabulates from the LLM alone. Confident and wrong. & Empty for this topic \\
\addlinespace
2 & \textbf{Unskilled, aware} & System 2 & Stings accumulating. The instructor said ``wrong foot.'' Agent searches effortfully, still fails often. Every response costs effort. & Sting-heavy, low quality\_scores \\
\addlinespace
3 & \textbf{Skilled, aware} & System 2 & Quality episodes with high scores. Agent retrieves reliably but runs the full recall pipeline every time. Counting one-two-three. & Rich, quality-scored, actively recalled \\
\addlinespace
4 & \textbf{Skilled, unaware} & System 1 & Episode promoted to hardcoded algorithm. Agent responds without searching. Hears the music and moves. & Episode promoted out --- became code \\
\bottomrule
\end{tabularx}
\caption{Four phases of competence mapped to Kahneman's System 1/System 2}
\end{table}

The transition from Phase 1 to Phase 2 is the most important moment in the system's life. It is the moment the agent stops being confidently wrong and starts being aware of its own ignorance. This transition is triggered by the first sting --- the first time a user says ``wrong, because X.'' That single sentence of explanation creates the first corrective neuron. The agent now knows it has a problem.

Phase 2 to Phase 3 is practice. Repeated interactions, repeated feedback, repeated retrieval. The quality scores rise. The stings get resolved --- their outcomes change from ``unresolved'' to ``resolved'' as corrective principles are confirmed. The agent gets the steps right, but it's still counting.

Phase 3 to Phase 4 is mastery. An episode that has been recalled hundreds of times with consistent positive ratings is a candidate for promotion from database retrieval to hardcoded algorithm. In our three-tier distillation model: Tier 3 (general LLM) teaches Tier 2 (agent's own LLM) teaches Tier 1 (hard algorithm in code). The episode becomes code. The agent stops searching and starts knowing. The counting stops. The music takes over.

\subsection{Brain-Like Agent Components}

Three memory capacities work together to support the four-phase progression:

\textbf{Hippocampus (short-term working memory)} holds the current interaction context. It detects when the current situation differs from the pattern the algorithm was promoted from. If the match is poor, the system escalates from System 1 back to System 2 --- from automatic response back to deliberate retrieval. This is the dancer noticing that the music has changed tempo.

\textbf{Retrospection (periodic measurement)} measures outcomes against expectations on a scheduled basis, even when every answer appears correct. The question is not ``is the answer wrong?'' but ``is there a better answer?'' This is the mechanism that prevents institutional calcification --- the slow death that comes from never questioning what works. It asks the question that no automated reward signal asks: could we be more right?

\textbf{Episodic memory (continued accumulation)} keeps recording new episodes even after promotion. If a promoted algorithm starts generating stings --- users saying ``wrong, because X'' about something the system thought it knew --- the knowledge cycles back from Phase 4 to Phase 2. The episode store catches the drift. The dancer goes back to class.

No memory without retrospection. No retrospection without measurement. No measurement without memory markers. The three form a closed loop. Break any link and the system stops learning.

\subsection{The Phase 4 Danger: Retrospection is Mandatory Behavior}

Phase 4 is powerful and dangerous. ``Skilled and unaware'' means the system no longer questions what works. Promoted algorithms feel correct because they have always worked --- but the environment changes. A dancer who never reviews their form develops habits that feel natural but limit growth. An institution that never reexamines its procedures becomes a bureaucracy. Phase 4 without retrospection is not expertise. It is habit.

What is the difference between knowledge and wisdom? Scars. Knowledge is what you have been told. Wisdom is what you have learned from getting it wrong, understanding why, and checking whether your correction still holds. Experience applies intelligence with greater skill --- but only if the experience is subjected to ongoing retrospection. The moment you stop checking is the moment your expertise begins to decay. This is as true for organizations as it is for dancers, and it is as true for AI systems as it is for organizations.

\subsection{The Neural Analogy}

The correspondence between our architecture and biological neural mechanisms is structural, not metaphorical:

\begin{table}[h]
\centering
\small
\begin{tabularx}{\textwidth}{>{\raggedright}p{3.5cm} >{\raggedright}p{3cm} >{\raggedright\arraybackslash}p{6.5cm}}
\toprule
\textbf{Neural Mechanism} & \textbf{System Component} & \textbf{Mechanism} \\
\midrule
Synapse & Episode & Discrete unit of stored experience connecting stimulus to response \\
\addlinespace
Synaptic strength & quality\_score & Scalar weight determining retrieval priority; adjusted by feedback \\
\addlinespace
Long-term potentiation (LTP) & Thumbs up & Repeated positive feedback increases quality\_score, strengthening the retrieval pathway \\
\addlinespace
Long-term depression (LTD) & Thumbs down & Negative feedback decreases quality\_score, weakening the pathway \\
\addlinespace
Hebbian learning & Co-recall promotion & Episodes retrieved together in successful responses are both promoted --- they ``fire together, wire together'' \\
\addlinespace
Pattern completion & Associative recall & Partial semantic match triggers full episode retrieval --- the system completes the pattern from a fragment \\
\addlinespace
Frequency-dependent plasticity & recall\_count & Episodes activated more frequently accumulate higher retrieval influence \\
\addlinespace
Synaptic pruning & Quality threshold & Episodes with consistently negative quality\_score fall below retrieval threshold --- functionally eliminated without deletion \\
\addlinespace
Cross-region connectivity & Cross-agent queries & Episodes from one agent's domain inform another agent's responses --- knowledge transfers across functional boundaries \\
\addlinespace
Consolidation & Phase 3$\rightarrow$4 promotion & Frequently recalled, high-quality episodes are promoted from database retrieval to hardcoded algorithm --- episodic memory becomes procedural memory \\
\bottomrule
\end{tabularx}
\caption{Structural correspondence between biological neural mechanisms and the episode architecture}
\end{table}

\subsection{Why This Is Not Metaphor}

Traditional AI training adjusts billions of model weights across the entire network. A training run modifies the global weight distribution to reduce aggregate loss. No individual weight has semantic meaning. The granularity of learning is the model.

Our system adjusts individual episode weights based on specific user feedback on specific interactions. Each quality\_score adjustment has direct semantic meaning: this specific episode, in this specific context, produced a response that the user rated positively or negatively for a stated reason. The granularity of learning is the individual memory.

This matches biological neural plasticity more closely than model training does. The brain does not retrain all 100 trillion synapses when you learn that the stove is hot. It strengthens the specific pathways associated with that specific experience. Our quality\_score adjustment operates at the same granularity.

The emergent property is the same: intelligence arises not from the sophistication of individual components but from the density and quality of connections between them. A neuron is a simple threshold function. An episode is a simple database record. But 100 billion weighted neurons produce consciousness, and a sufficiently dense network of quality-scored episodes produces domain expertise --- without any individual component being ``intelligent.''

%% ====================================================================
\section{Related Work}

\subsection{Agent Memory Architectures}

Mem0 \cite{mem0} provides production middleware for agent memory with remember, recall, and forget operations, establishing the first broad comparison of memory approaches on the LoCoMo benchmark (ECAI 2025). Zep \cite{zep} builds on the Graphiti temporal knowledge graph to approximate episodic-to-semantic consolidation for conversational contexts. MEMTIER \cite{memtier} proposes a tiered memory architecture for long-running autonomous agents, analyzing retrieval bottlenecks across memory layers. Redis \cite{redis} documents long-term memory architectures using vector similarity search with recency and importance scoring.

These systems treat memory as a retrieval optimization problem. None incorporate quality scoring from human feedback on individual memories.

\subsection{Reinforcement Learning on Agent Memory}

MemRL \cite{memrl} is the closest precursor to our work. It employs a Two-Phase Retrieval mechanism that filters candidates by semantic relevance and then selects them based on learned Q-values, with utilities refined via environmental feedback in a trial-and-error manner. The critical difference: MemRL's learning signal comes from automated task success or failure. Our signal comes from human explanation of why the answer was wrong.

RoMeRL \cite{romerl} addresses the memory-reward trap in self-evolving agent memory, balancing feedback coverage with utility. Memory-R1 \cite{memoryr1} applies reinforcement learning to memory management decisions. The Memory for Autonomous LLM Agents survey \cite{memsurvey} provides a comprehensive taxonomy of mechanisms including consolidation, intelligent forgetting, and conflict resolution.

All of these systems use automated reward signals. None require the user to explain the failure. The explanation --- the ``why'' --- is where the principle lives.

\subsection{Multi-Agent Shared Memory}

Governed Shared Memory \cite{govmem} (ICML 2026) introduces governance rules for multi-agent memory pools, addressing access control and conflict resolution. Multi-Agent Transactive Memory \cite{transactive} (ICLR 2026) models agent specialization and knowledge sharing, drawing on organizational psychology.

These approaches use shared memory pools --- a common store that all agents access. Our architecture uses per-agent databases with cross-agent queries: each agent owns its experience, and sharing is a deliberate act of querying another agent's store. This preserves sovereignty --- each agent's rated episodes reflect its own domain experience --- while enabling the network effect of cross-domain knowledge transfer.

The gap in the literature: no existing system combines (a) per-agent databases, (b) cross-agent episode queries, (c) human-sourced quality feedback on individual episodes, and (d) a competence-phase model that includes promotion to hardcoded algorithms with mandatory retrospection.

%% ====================================================================
\section{Architecture}

\subsection{Episode Structure}

An episode is a structured record of a significant event. Unlike log entries (unstructured text) or memory embeddings (vector-only), episodes carry both structured metadata for precise querying and vector embeddings for similarity search.

\textbf{Schema (PostgreSQL):}

\begin{lstlisting}
episode_id       VARCHAR(20) UNIQUE    -- EP-xxxxxxxxxxxx
dt_created       BIGINT                -- milliseconds since epoch
dt_start         BIGINT                -- when the episode began
dt_end           BIGINT                -- when it resolved (if resolved)
episode_type     VARCHAR(30)           -- tfts|fault|scar|commerce|
                                       -- session|pattern|sting
domain           VARCHAR(20)           -- agent's operational domain
title            TEXT                   -- one-line summary
narrative        TEXT                   -- full story: what happened,
                                       -- tried, learned
principle        TEXT                   -- the lesson extracted (if any)
actors           JSONB                  -- who was involved
outcome          VARCHAR(20)           -- resolved|unresolved|ongoing
severity         VARCHAR(20)           -- lesson|scar|win
quality_score    FLOAT (in metadata)   -- user-feedback-adjusted (EMA)
recall_count     INTEGER               -- how often surfaced
last_recalled    BIGINT                -- last retrieval timestamp
\end{lstlisting}

\textbf{Vector store (ChromaDB):} Each episode's title, narrative, and principle are concatenated and embedded in a cosine-similarity collection. Metadata (episode\_type, domain, outcome, severity) enables filtered search.

\textbf{Episode types reflect the source of the knowledge:}
\begin{itemize}
    \item \textbf{tfts} (try-fail-try-succeed): A complete problem-solving arc with multiple attempts, revealing the principle that made the final attempt obvious in retrospect.
    \item \textbf{scar}: An experience that cost something --- time, data, trust. A scar episode records the cost alongside the lesson. These are the episodes that carry wisdom.
    \item \textbf{sting}: A corrective episode created from a user's thumbs-down explanation. Unresolved by default --- the sting records the failure but not yet the fix.
    \item \textbf{fault}, \textbf{commerce}, \textbf{session}, \textbf{pattern}: Domain-specific event types for mechanical faults, commercial transactions, development sessions, and recurring patterns respectively.
\end{itemize}

\subsection{Associative Recall}

Every agent interaction triggers associative recall before the agent responds:

\begin{enumerate}
    \item \textbf{Vector search:} Query the episodes collection with the current question or event text. Retrieve top-$k$ candidates (default $k=3$) by cosine similarity.

    \item \textbf{Quality-adjusted ranking:} Raw vector distance is adjusted by the episode's quality score:
    \begin{equation}
        \text{effective\_distance} = \text{vector\_distance} - (\text{quality\_score} \times 0.1)
    \end{equation}
    Episodes with positive quality scores (validated by user feedback) surface higher. Episodes with negative quality scores (associated with bad responses) sink. This is the Hebbian mechanism: connections that have proven useful are strengthened; connections associated with failure are weakened.

    \item \textbf{Threshold filtering:} Only episodes below the associative recall threshold (default 0.50 cosine distance, applied to effective\_distance) are surfaced. This prevents noise --- irrelevant episodes with good quality scores cannot override semantic dissimilarity.

    \item \textbf{Reinforcement:} Every recalled episode has its recall\_count incremented and last\_recalled updated. Frequently recalled episodes accumulate higher recall\_count values, enabling future analysis of which episodes are load-bearing knowledge and which are dormant.
\end{enumerate}

\subsection{The Small-Stings Feedback Loop}

The feedback mechanism is deliberately asymmetric:

\textbf{Thumbs up} requires no explanation. The quality\_score of all episodes surfaced in the response increases. The cost to the user is zero. Positive reinforcement flows freely.

\textbf{Thumbs down} requires one sentence: why. This constraint is the core innovation. The explanation serves three purposes:

\begin{enumerate}
    \item \textbf{It creates a new episode.} The sting episode records the failure and the user's diagnosis. It is immediately indexed in the vector store.

    \item \textbf{It demotes the episodes that were surfaced.} Their quality scores decrease. Future retrieval will rank them lower, or if quality is sufficiently negative, they fall below the recall threshold --- functionally pruned.

    \item \textbf{It surfaces on similar future queries.} Because the sting is vector-indexed with the content of the failed interaction, it will be retrieved when a similar question arrives. The agent now has a corrective memory: ``Last time this type of question came up, the answer was wrong because X.''
\end{enumerate}

\textbf{Quality score update} uses an exponential moving average:
\begin{equation}
    \text{quality\_new} = \text{quality\_old} \times 0.8 + \text{grade\_delta} \times 0.2
\end{equation}
Where grade\_delta is $+1.0$ for thumbs up and $-1.0$ for thumbs down. The 0.8 decay preserves the influence of historical ratings while giving recent feedback sufficient weight to shift behavior. An episode must accumulate multiple consecutive negative ratings before being effectively pruned, preventing single outlier ratings from destroying useful knowledge.

\subsection{Multi-Agent Database Architecture}

Each agent has its own PostgreSQL database and ChromaDB vector store. The schema is identical across all agents --- created by a single setup script --- but the data is specific to each agent's domain.

\begin{table}[h]
\centering
\small
\begin{tabular}{lllll}
\toprule
\textbf{Agent} & \textbf{Domain} & \textbf{Database} & \textbf{Log Type} & \textbf{Specialty} \\
\midrule
Alice & Commerce & agent\_alice & commerce & Pricing, billing, data quality \\
Andi & Production & agent\_andi & commerce & WCHQ server, user-facing \\
Noelle & Design & agent\_noelle & validation & Network validation, load balancing \\
Natalie & Routing & agent\_natalie & routing & Trip plans, route optimization \\
Nora & Telemetry & agent\_nora & telemetry & Vehicle sensors, navigation \\
Sally & Station & agent\_sally & station & Slot registry, parking queue \\
Allie & Cross-domain & allie & general & Coordination, synthesis \\
\bottomrule
\end{tabular}
\caption{Seven deployed agents with per-agent databases}
\end{table}

\textbf{Core tables} (universal, every agent): episodes, observations, agent\_messages, agent\_facets, sessions, vector\_index, tfts.

\textbf{Agent-specific log table} (custom columns per domain): Alice logs model\_name, record\_id, customer\_id. Nora logs pod\_name, sensor, value\_raw, value\_calibrated.

\textbf{Cross-agent queries} are enabled through the agent message bus. When an agent encounters a question that spans domains, it can query another agent's episode store. The querying agent rates the relevance of cross-agent episodes it receives, contributing quality signal to the source agent's episodes. This creates a network effect: each agent's use of another agent's episodes improves the quality scoring for both.

\subsection{Escalation with Episode Context}

The system employs a three-tier escalation chain:

\textbf{Tier 1:} Local LLM (Ollama, running on the agent's hardware). Fast, private, free. Always first.

\textbf{Tier 2:} WCHQ Alice (shared LLM at WebClerk headquarters). Better model, \$4/person/month. Triggered when local confidence falls below 0.40.

\textbf{Tier 3:} WCHQ Claude (Anthropic's Claude, managed centrally). Best reasoning, \$9/person/month. Triggered when WCHQ Alice is also low-confidence.

The critical design decision: \textbf{escalation payloads include recalled episodes.} The upstream model receives not just the question but the relevant episodes --- the team's accumulated experience on that topic. Claude doesn't reason from scratch. It reasons with context that includes prior failures, their explanations, and the principles extracted from them. The response flows back and is stored as a new episode at the local agent, enriching the episode store for future local retrieval.

This creates a distillation dynamic: expensive Tier 3 reasoning is captured as an episode that enables future Tier 1 retrieval. Over time, the system escalates less because the episode store grows richer. The dance floor produces its own instructors.

%% ====================================================================
\section{Metcalfe's Law Applied to Agent Intelligence}

Value is proportional to connections squared. More agents means more episodes. More episodes means more cross-agent queries. More queries means more ratings. More ratings means better retrieval for everyone.

Each new agent makes every existing agent smarter. Small packets of experience --- episodes --- flow across the network, rated by each agent that uses them, creating compound intelligence that no single agent could develop alone. This is in direct contrast to the centralized training paradigm, where intelligence flows from center to edge. In our architecture, intelligence flows from edge to edge, rated and reinforced at every connection point.

We offer every new installation two months of full access at no charge. This is not marketing. It is growing neural density. Every real user with real data generates episodes that no synthetic training data can replicate. Their stings teach the system things no prompt engineer would think to include. The trial period is the 200 hours on the dance floor.

%% ====================================================================
\section{Implementation and Deployment}

\subsection{System Overview}

The architecture is deployed across four operational domains:

\textbf{Commerce (WebClerk):} Alice and Andi handle pricing queries, data quality validation, customer pattern detection, and onboarding guidance. Alice runs locally at each installation. Andi runs at the WCHQ production server, handling escalations.

\textbf{Transit network design (JPods Physical Internet\textsuperscript{\textregistered}):} Noelle validates station capacity, guideway connections, and network topology. Her episodes record build faults, validation outcomes, and design decisions.

\textbf{Vehicle operations (JPods routing and vehicles):} Natalie plans routes and dispatch sequences. Nora processes sensor telemetry from vehicle hardware. Sally manages station slot assignment and parking queues.

\textbf{Cross-domain coordination (Allie):} Allie synthesizes across all domains, maintaining the cross-domain episode store and running nightly retrospection.

The system was seeded with 193 initial episodes: 120 try-fail-try-succeed (TFTS) arcs extracted from development session records, and 73 operational scars from the team's experience database.

\subsection{Technology Stack}

\begin{itemize}
    \item \textbf{PostgreSQL} for structured episode storage and agent-specific log tables
    \item \textbf{ChromaDB} (PersistentClient, cosine similarity) for vector embeddings
    \item \textbf{Ollama} (gpt-oss:20b base model) for local LLM inference
    \item \textbf{MCP (Model Context Protocol)} for agent tool exposure to Claude Code
    \item \textbf{Python} for all agent servers, ingestion scripts, and coordination logic
\end{itemize}

All components run on local hardware. No cloud dependency for core operations. Escalation to WCHQ is optional and used only when local confidence is insufficient.

\subsection{Episode Ingestion Pipeline}

Episodes enter the system through four channels:

\begin{enumerate}
    \item \textbf{Automated ingestion:} A script parses existing TFTS markdown files and retro.db experience records into structured episodes, storing them in both PostgreSQL and ChromaDB.

    \item \textbf{Session creation:} During development sessions, the operator creates episodes for significant decisions, faults, and debugging arcs using the alice\_episode\_create MCP tool.

    \item \textbf{Sting creation:} When a user gives a thumbs-down rating with an explanation, the system automatically creates a sting episode and indexes it.

    \item \textbf{Cross-agent sharing:} Episodes are shared between agents via the agent\_messages table. The receiving agent indexes the episode in its own ChromaDB collection.
\end{enumerate}

\subsection{Operational Metrics}

Five metrics monitor system health:

\begin{itemize}
    \item \textbf{Episode creation rate:} New episodes per day, by type and domain. Stagnation indicates the feedback loop has broken.
    \item \textbf{Recall hit rate:} Percentage of interactions where associative recall surfaces at least one episode below threshold.
    \item \textbf{Quality score distribution:} A healthy system trends positive over time.
    \item \textbf{Sting resolution rate:} Percentage of stings that subsequently acquire a principle. Low resolution rate indicates stings are being created but not learned from.
    \item \textbf{Cross-agent query volume:} Growth indicates the network effect is active.
\end{itemize}

%% ====================================================================
\section{Comparison with Existing Approaches}

\begin{table}[h]
\centering
\small
\begin{tabularx}{\textwidth}{>{\raggedright}p{2.2cm} >{\raggedright}p{2.2cm} >{\raggedright}p{2.2cm} >{\raggedright}p{2.2cm} >{\raggedright\arraybackslash}p{3cm}}
\toprule
\textbf{Feature} & \textbf{MemRL} & \textbf{Mem0} & \textbf{Gov.\ Shared Memory} & \textbf{This work} \\
\midrule
Quality signal source & Automated task reward & None & Governance rules & Human feedback with required explanation \\
\addlinespace
Granularity & Q-values per memory & Per-entry importance & Per-document & Per-episode quality\_score \\
\addlinespace
Learning mechanism & RL optimization & Recency / importance & Rule-based access & Hebbian (EMA from user grades) \\
\addlinespace
Negative feedback & Task failure (binary) & N/A & Access denial & Explanation required $\rightarrow$ sting episode created \\
\addlinespace
Multi-agent & Single agent & Single agent & Shared pool & Per-agent databases + cross-agent queries \\
\addlinespace
Scaling mechanism & Retrain policy & Add data & Add agents to pool & More use = more rated episodes (no retraining) \\
\addlinespace
Competence phases & Not modeled & Not modeled & Not modeled & Four phases with Phase 4 retrospection \\
\addlinespace
Neural analogy & Implicit (Q-values $\approx$ synaptic weights) & None & None & Explicit (full Hebbian correspondence) \\
\addlinespace
Model dependency & Requires capable model & Model-agnostic & Model-agnostic & Model-agnostic (database is the intelligence) \\
\addlinespace
Data sovereignty & Cloud & Cloud & Cloud & Local-first (agent owns its database) \\
\bottomrule
\end{tabularx}
\caption{Comparison with existing agent memory approaches}
\end{table}

The most significant differentiator is the quality signal source. MemRL learns from whether the task succeeded or failed --- a binary signal with no explanatory content. Our system learns from the user's explanation of why the response was wrong. The explanation contains the principle. The principle is what prevents the next failure. No automated reward function generates principles.

The sting mechanism is a secondary differentiator. In MemRL, a failed task reduces the Q-value of the retrieved memory. In our system, a failed response creates a new corrective memory that actively surfaces on similar future queries. The failure doesn't just weaken the wrong path --- it creates the right path. This is the difference between long-term depression alone (weakening a synapse) and long-term depression paired with compensatory potentiation (weakening one pathway while strengthening a corrective one).

%% ====================================================================
\section{Limitations and Future Work}

\textbf{Local LLM reasoning capacity.} The local language model (gpt-oss:20b) has limited ability to synthesize across multiple recalled episodes. It functions primarily as a retrieval presentation layer. As local models improve, Phase 3 performance will improve proportionally --- without any changes to the episode architecture.

\textbf{Heuristic EMA parameters.} The quality score update parameters (decay 0.8, weight 0.2) and quality-distance factor (0.1) were chosen based on engineering judgment. A systematic study of parameter sensitivity would improve generalizability.

\textbf{Manual cross-agent routing.} Cross-agent queries currently require explicit routing. Automatic domain detection based on question content would increase the network effect without operator intervention.

\textbf{No explicit forgetting mechanism.} Episodes are functionally pruned by negative quality scores but never deleted. A consolidation mechanism --- similar to biological memory consolidation from hippocampus to neocortex --- would compress validated episodes into semantic summaries while archiving originals.

\textbf{Longitudinal validation needed.} A rigorous study measuring recall precision, sting resolution rate, quality score trajectory, and cross-agent query patterns over 6--12 months would provide empirical validation of the learning dynamics described here.

\textbf{Phase 4 promotion criteria.} The criteria for promoting an episode from database retrieval to hardcoded algorithm are not yet formalized. Automated promotion with confidence thresholds and mandatory retrospection scheduling would complete the four-phase lifecycle.

%% ====================================================================
\section{Conclusion}

What is the difference between knowledge and wisdom? Scars. A language model has knowledge --- vast, general, and reset every session. Wisdom requires scars: specific memories of specific failures with specific explanations of why they occurred. Experience applies intelligence with greater skill.

We have presented an architecture where AI agent intelligence accumulates in quality-scored database episodes rather than in model parameters. The architecture is structurally analogous to biological neural systems: episodes function as synapses, user feedback functions as long-term potentiation and depression, cross-agent queries function as cross-region connectivity, and the four-phase competence model traces the lifecycle of knowledge from conscious incompetence through to unconscious competence with mandatory retrospection.

The moat for AI-augmented business systems is not the language model --- models are commodities that converge. The moat is the accumulated rated experience specific to a domain, its customers, and its operations. That database of quality-scored episodes cannot be replicated by training a larger model. It can only be earned through use: through practice on the dance floor, through stings from the missteps, through the 200 hours that build the neurons no lecture hall can provide.

Intelligence combined with experience based skills results from accumulated rated experience across connected agents. The database is the brain. The language model is the mouth. The user's feedback --- especially the small sting of explaining why an answer was wrong --- is the learning signal. The system scales with use and belongs to the user, not to the platform.

Learn to dance, it will help you program better ;-)

%% ====================================================================
\bibliographystyle{plain}

\begin{thebibliography}{12}

\bibitem{mem0}
Mem0. ``Long-Term Memory for AI Agents.'' ECAI 2025. \url{https://mem0.ai}

\bibitem{zep}
Zep / Graphiti. Temporal knowledge graph for agent memory. \url{https://getzep.com}

\bibitem{memtier}
MEMTIER: ``Tiered Memory Architecture and Retrieval Bottleneck Analysis for Long-Running Autonomous AI Agents.'' arXiv:2605.03675, 2025.

\bibitem{redis}
Redis. ``Long-Term Memory Architectures for AI Agents.'' redis.io/blog, 2025.

\bibitem{memrl}
MemRL: ``Self-Evolving Agents via Runtime Reinforcement Learning on Episodic Memory.'' arXiv:2601.03192, 2026.

\bibitem{romerl}
RoMeRL: ``Balancing Feedback Coverage and the Memory-Reward Trap in Self-Evolving Agent Memory via Reduced-Order Utility States.'' arXiv:2608.02508, 2026.

\bibitem{memoryr1}
Memory-R1: ``Enhancing Large Language Model Agents to Manage and Retrieve Memory.'' arXiv:2508.19828, 2025.

\bibitem{memsurvey}
``Memory for Autonomous LLM Agents: Mechanisms, Evaluation, and Emerging Frontiers.'' arXiv:2603.07670, 2026.

\bibitem{govmem}
``Governed Shared Memory for Multi-Agent LLM Systems.'' arXiv:2606.24535, ICML 2026.

\bibitem{transactive}
``Multi-Agent Transactive Memory.'' arXiv:2606.19911, ICLR 2026.

\bibitem{kahneman2011}
Kahneman, D. \textit{Thinking, Fast and Slow.} Farrar, Straus and Giroux, 2011.

\bibitem{hebb1949}
Hebb, D.O. \textit{The Organization of Behavior.} Wiley, 1949.

\end{thebibliography}

\end{document}
